\section{MARL Algorithms with Theory}\label{sec:algorithms}
This section  provides a selective  review of MARL algorithms, and categorizes them according to the tasks to address.  Exclusively, we  review here the works that are focused on the theoretical studies only, which are mostly built upon the two representative MARL frameworks, fully observed Markov games and extensive-form games, introduced in \S\ref{subsec:MARL_framework}. A brief summary  on MARL for partially observed Markov games in  \emph{cooperative}  settings, namely, the Dec-POMDP problems, is also provided below in \S\ref{subsec:cooperative_MARL_alg},  due to their relatively  more mature   theory than that of MARL for general partially observed Markov games.
 
\subsection{Cooperative Setting}\label{subsec:cooperative_MARL_alg}

Cooperative MARL constitutes a great portion of  MARL settings, where all agents collaborate with each other to achieve some shared goal. Most cooperative MARL algorithms backed by  theoretical analysis are devised for the following more specific settings. 

%which are mostly built upon the Markov games model, possibly with partiall . 

\subsubsection{Homogeneous Agents}\label{subsubsec:homo_agents}

A majority of cooperative MARL settings  involve \emph{homogeneous} agents with a \emph{common} reward function that aligns all agents' interests. In the extreme case with large populations of agents, such a homogeneity also indicates that the agents play an \emph{interchangeable}  role in the system evolution, and can hardly be distinguished from each other. We elaborate more on  homogeneity below. 

\vspace{6pt}
\noindent{\bf Multi-Agent MDP \& Markov Teams}
\vspace{3pt}

Consider a Markov game as in Definition \ref{def:Markov_Game} with $R^1=R^2=\cdots=R^N=R$, where  the reward $R:\cS\times\cA\times\cS\to\RR$ is influenced by the joint action in $\cA=\cA^1\times\cdots\times \cA^N$. 
%We note that such a model is also referred to as \emph{multi-agent MDPs} in the  AI community \cite{boutilier1996planning,lauer2000algorithm}, and \emph{Markov/dynamic  teams} in the game/control theory community \cite{ho1980team,wang2003reinforcement,ma}. 
As a result, the Q-function is identical for  all agents. 
%, and more intriguingly,  finding the optimal joint policy reduces  to finding the NE policy with the highest  return.  
%In other words, from the team's perspective, this problem reduces to a single-agent RL one, if all agents coordinate properly. 
Hence, a straightforward algorithm proceeds by performing the standard 
Q-learning update \eqref{equ:Q_learning}  at each agent, but taking the $\max$ over the joint action space $a'\in\cA$. Convergence to the optimal/equilibrium Q-function has been established  in \cite{szepesvari1999unified,littman2001value}, when both state and action spaces are finite.

However,  convergence of the Q-function does not necessarily imply that of the equilibrium policy  for the Markov team, as any combination of  equilibrium policies extracted at each agent may not constitute an equilibrium policy, if the equilibrium policies are non-unique, and the agents fail to agree on which one to select. Hence, convergence to the NE policy is only guaranteed if either the equilibrium is assumed to be unique \cite{littman2001value}, or the agents are coordinated for equilibrium selection.  The latter idea has first been validated in the cooperative repeated games setting \cite{claus1998dynamics},  a special case of Markov teams with a singleton state, where  the agents are joint-action learners (JAL),  maintaining a Q-value for joint actions, and learning  empirical models of all others. Convergence to equilibrium point is claimed  in \cite{claus1998dynamics}, without a formal proof. 
For the actual Markov teams, this coordination has been exploited in  \cite{wang2003reinforcement}, which proposes \emph{optimal adaptive learning} (OAL), the first MARL algorithm   with provable  convergence to the equilibrium policy.  Specifically, OAL first learns the game structure, and constructs virtual games at each state that are \emph{weakly acyclic} with respect  to (w.r.t.) a biased set. OAL can be shown to converge to the NE, by introducing the biased adaptive play learning algorithm for the constructed  weakly acyclic games, motivated from \cite{young1993evolution}.  


  
Apart from equilibrium selection, another subtlety   special to Markov teams (compared to single-agent RL) is the necessity to address the scalability issue, see  \S\ref{subsec:scala_issue}. As   independent Q-learning may fail to converge \cite{tan1993multi}, one early attempt toward developing  scalable while convergent algorithms for MMDPs is \cite{lauer2000algorithm}, which advocates a distributed Q-learning algorithm  that converges for \emph{deterministic} finite  MMDPs. Each agent maintains only a Q-table of state $s$ and local action $a^i$, and successively takes maximization over the joint action $a'$. No other agent's actions and their histories can be acquired by each individual agent. 
Several other heuristics (with no theoretical backing) regarding either reward or value function factorization  have been proposed to mitigate the scalability issue \cite{guestrin2002coordinated,guestrin2002multiagent,kok2004sparse,sunehag2018value,rashid2018qmix}.  Very recently,  \cite{pmlr-v97-son19a} provides a rigorous characterization of  conditions that justify  this value factorization idea. Another recent theoretical work  along this direction is \cite{qu2019exploiting},    which imposes a special dependence structure, i.e., a one-directional tree, so that the (near-)optimal policy of the overall MMDP can be provably well-approximated  by \emph{local policies}. 
More recently, 
\cite{yongacoglu2019learning}  has studied    common interest games, which  includes Markov teams as an example, and develops a \emph{decentralized} RL algorithm that relies on only states,  local   actions and rewards. With the same information structure as independent Q-learning \cite{tan1993multi}, the  algorithm is guaranteed to converge to \emph{team optimal} equilibrium policies, and not just equilibrium policies. This is important as in general, a suboptimal equilibrium can perform arbitrarily worse than
the optimal equilibrium  \cite{yongacoglu2019learning}. 

%, which take  only local states as input. This approximation in some sense   addresses both the scalability and partial observability issues mentioned  in \S\ref{sec:challenges} in some sense.   


For 
policy-based methods,  
to  our knowledge, the only convergence guarantee  for this setting exists in  \cite{perolat2018actor}.  The authors propose two-timescale actor-critic \emph{fictitious play} algorithms, where at the slower timescale, the actor mixes the current policy and the best-response one w.r.t. the local Q-value estimate, while at the faster timescale the critic performs policy evaluation, as if all agents' policies are stationary. Convergence is established for  \emph{simultaneous move
multistage games} with a common (also zero-sum, see \S\ref{sec:policy_comp}) reward, a special Markov team  with initial and absorbing states, and each state being visited only once. 
   
%Factorization-based algs: 
%Though independent Q-learning does not work, distributed Q-learning works for deterministic MMDPs \cite{lauer2000algorithm}. 
%Then mention other factoizations, but they are heuristics. Except recently,   guannan's recent work... in the track of "coordinate graph" \& handle scalability issues, with value factorization.



\vspace{6pt}
\noindent{\bf Markov Potential Games}
\vspace{3pt}

From a game-theoretic perspective, a more general framework to embrace cooperation is \emph{potential games} \cite{monderer1996potential}, where there exits some \emph{potential} function shared by all agents, such that   if any agent changes its policy unilaterally, the change in its reward equals (or proportions to) that in the potential function. Though most potential games are stateless,  an extension named \emph{Markov potential games} (MPGs)  has gained increasing attention for modeling \emph{sequential} decision-making \cite{gonzalez2013discrete,zazo2016dynamic}, which includes Markovian states whose evolution is affected   by the joint actions. 
Indeed, MMDPs/Markov teams constitute   a particular case of MPGs, with the potential function being the common reward; such dynamic games can also be viewed as being {\em strategically equivalent} to Markov teams, using the terminology in, e.g., \cite[Chapter $1$]{basar2018handbook}. Under this model, \cite{valcarcel2018learning} provides  verifiable  conditions for a Markov game to be an MPG, and shows the equivalence between  finding closed-loop  NE in MPG and solving a single-agent optimal control problem. Hence,  single-agent RL algorithms are then enabled to solve this MARL problem. 
 

\vspace{6pt}
\noindent{\bf Mean-Field Regime}
\vspace{3pt}

Another  idea toward tackling the scalability issue  is to take the setting to the \emph{mean-field} regime, with an extremely large number of homogeneous agents. 
Each agent's effect on the overall multi-agent system can thus become  infinitesimal, resulting in all agents being  interchangeable/indistinguishable. The interaction with other agents,  however, is captured simply by some mean-field quantity, e.g., the average state, or the empirical distribution of states. Each agent only needs to find the best response to the mean-field, which considerably simplifies the analysis.  
This mean-field view of multi-agent systems has been approached by the mean-field games (MFGs) model \cite{huang2003individual,huang2006large,lasry2007mean,bensoussan2013mean,tembine2013risk}, the team model with mean-field sharing \cite{arabneydi2014team,arabneydi2017new}, and the game model with mean-field actions \cite{pmlr-v80-yang18d}.\footnote{The difference between mean-field teams and mean-field games is mainly the solution concept: optimum versus equilibrium, as the difference between general dynamic team theory \cite{witsenhausen1971separation,yoshikawa1978decomposition,yuksel2013stochastic} and game theory \cite{shapley1953stochastic,filar2012competitive}. Although the former can be viewed as a special case of the latter, related works are usually reviewed separately in the literature.  We  follow here this convention. 
} 

MARL in these models have not been explored until recently, mostly in the non-cooperative setting of MFGs, see \S\ref{subsec:mixed_setting} for a more detailed review. Regarding the cooperative setting, recent work  \cite{subramanian2018reinforcement} studies RL for Markov teams with mean-field sharing \cite{arabneydi2014team,arabneydi2016linear,arabneydi2017new}. Compared to MFG, the model considers agents that share a common reward function depending only on the local state and the mean-field, which encourages cooperation among the  agents. Also, the term mean-field refers to the \emph{empirical average} for the states of \emph{finite} population, in contrast to the \emph{expectation} and \emph{probability distribution} of \emph{infinite} population in MFGs. 
 Based on the dynamic programming decomposition for the specified model  \cite{arabneydi2014team}, several popular  RL algorithms are easily translated to address this setting \cite{subramanian2018reinforcement}. More recently, \cite{carmona2019linear,carmona2019model}  approach the problem from a mean-field control (MFC) model,  to model large-population  of cooperative decision-makers. Policy gradient methods are proved to converge for linear quadratic MFCs in \cite{carmona2019linear}, and mean-field Q-learning is then shown to converge for general MFCs \cite{carmona2019model}.
 
%Mean-field teams formulation, LQ setting \cite{arabneydi2016linear}, general setting with mean-field sharing \cite{arabneydi2014team,arabneydi2017new}. Then mention the only work that considers RL in mean-field teams \cite{subramanian2018reinforcement},  and preview MFG. 
%\issue{also mention the ``mean-field'' regime.}
%{The idea of Mean-field teams but if there any learning work? what about Mahajan's student's thesis?}


 
\subsubsection{Decentralized Paradigm with Networked Agents}\label{subsubsec:networked_paradigm}
 
Cooperative agents in  numerous  
practical multi-agent systems are not always  homogeneous. Agents may have different preferences, i.e., reward functions, while  they still form a team to maximize the return of the \emph{team-average} reward $\bar R$, with $\bar R(s,a,s')=N^{-1}\cdot\sum_{i\in\cN}R^i(s,a,s')$. More subtly, the reward function is sometimes not sharable with others, as the preference is kept private to each agent.  
This setting finds broad applications in  engineering systems as sensor networks  \cite{rabbat2004distributed},  smart grid \cite{dall2013distributed,zhang2018dynamicpower}, intelligent transportation systems \cite{adler2002cooperative,zhang2018dynamic}, and robotics \cite{corke2005networked}. 


Covering the homogeneous  setting in \S\ref{subsubsec:homo_agents} as a special case, the specified one definitely  requires more coordination, as, for example, the  global value function  cannot be estimated locally without knowing other agents' reward functions. With a central controller, most MARL algorithms reviewed in   \S\ref{subsubsec:homo_agents}  directly apply, since the controller can collect and average the rewards, and distributes the information to all agents. Nonetheless,  such a controller may not exist in 
most aforementioned applications,  due to  either cost, scalability, or robustness concerns \cite{rabbat2004distributed,dall2013distributed,zhang2018distributedauto}. 
Instead,   the agents may be able to share/exchange information with their neighbors over  a possibly time-varying and sparse communication network, as  illustrated in    Figure \ref{fig:info_struc} (b). 
Though MARL under this \emph{decentralized/distributed}\footnote{Note that hereafter  we use \emph{decentralized} and \emph{distributed} interchangeably for describing this paradigm. } paradigm  is   imperative, it is  relatively less-investigated, in comparison to the extensive   results on distributed/consensus algorithms  that solve \emph{static/one-stage} optimization problems \cite{nedic2009distributed,agarwal2011distributed,jakovetic2011cooperative,tu2012diffusion}, which, unlike RL, involves no system \emph{dynamics}, and does not maximize the \emph{long-term}  objective as a \emph{sequential-decision making} problem. 

\vspace{6pt}
\noindent{\bf Learning Optimal Policy}
\vspace{3pt}

The most significant goal is to learn the optimal joint policy, while each agent only accesses to  local and neighboring information over the network.  
The idea of MARL with networked agents dates back to \cite{varshavskaya2009efficient}. 
To our knowledge, 
 the first provably convergent MARL algorithm under this setting is  due to \cite{kar2013cal}, which incorporates  the idea of  \emph{consensus $+$ innovation} to the standard Q-learning algorithm, yielding the \emph{$\mathcal{QD}$-learning} algorithm with the following update
\$
Q^i_{t+1}(s,a)\leftarrow & Q^i_{t}(s,a)+\alpha_{t,s,a}\Big[R^i(s,a)+\gamma\max_{a'\in\cA}Q^i_t(s',a')-Q^i_t(s,a)\Big] \notag \\
& \qquad -\beta_{t,s,a}\sum_{j\in\cN^i_t}\big[Q^i_t(s,a)-Q^j_t(s,a)\big],
\$
where $\alpha_{t,s,a},~\beta_{t,s,a}>0$ are  stepsizes, $\cN^i_t$ denotes the set of neighboring agents of agent $i$, at time $t$.  Compared to the Q-learning update \eqref{equ:Q_learning}, \emph{$\mathcal{QD}$}-learning appends an innovation term that captures the difference of Q-value estimates from its neighbors. With certain conditions on the stepsizes, the algorithm is guaranteed to converge to the optimum Q-function for the tabular setting.  

Due to the scalability issue, function approximation is vital in MARL, which necessitates  the development of policy-based  algorithms. Our work  \cite{zhang2018fully} proposes actor-critic algorithms for this setting. Particularly, each agent parameterizes its own policy $\pi^i_{\theta^i}:\cS\to\Delta(\cA^i)$ by some parameter $\theta^i\in\RR^{m^i}$, the policy gradient of the return  is first derived as 
%\small
\#\label{eq:policy_gradient_thm}
		\nabla _{\theta^{i} } J({\theta}) &= \EE \left [ \nabla _{\theta^i} \log \pi _{\theta^i} ^i(s,a^i)\cdot  Q_{{\theta}}(s,a) \right ]
%		=\EE  \left [ \nabla _{\theta^i} \log \pi _{\theta^i}^i (s,a^i)\cdot  A_{{\theta}}(s,a) \right ],
%		\nonumber\\
%		&=\EE _{s \sim d_{{\theta}}, a \sim \pi_{\theta} } \left [ \nabla _{\theta^i} \log \pi _{\theta^i}^i (s,a^i)\cdot  A^i_{{\theta}}(s,a) \right ].
\#
\normalsize
where  $Q_{{\theta}}$ is the global value function corresponding to $\bar R$ under the joint policy $\pi_\theta$ that is defined as $\pi_\theta(a\given s):=\prod_{i\in\cN}\pi^i_{\theta^i}(a^i\given s)$. As an analogy to \eqref{equ:policy_grad}, the policy gradient in \eqref{eq:policy_gradient_thm}  involves the expectation of the product between  the local score function $\nabla _{\theta^i} \log \pi _{\theta^i} ^i(s,a^i)$ and the global  Q-function $Q_{{\theta}}$.  The latter, nonetheless, cannot be estimated individually at each agent. As a result, by parameterizing each local copy of $Q_{{\theta}}(\cdot,\cdot)$ as $Q_{{\theta}}(\cdot,\cdot;\omega^i)$ for agent $i$, we propose the following consensus-based TD learning for the critic step, i.e., for estimating $Q_{{\theta}}(\cdot,\cdot)$:
\#
\tilde{\omega}^i_{t}=\omega^i_{t}+\beta_{\omega,t}\cdot \delta^i_{t}\cdot \nabla_{\omega} Q_t(\omega^i_t),\qquad\qquad{\omega}^i_{t+1} = \sum_{j\in\cN}c_t(i,j)\cdot \tilde{\omega}^j_{t}
\label{equ:MARL_critic_1},
\#
where $\beta_{\omega,t}>0$ is the stepsize,  
$\delta^i_{t}$ is the local TD-error calculated using $Q_{{\theta}}(\cdot,\cdot;\omega^i)$. The first relation in \eqref{equ:MARL_critic_1} performs the  standard TD update, followed by a weighted  combination of the neighbors' estimates $\tilde{\omega}^j_{t}$. The weights here, $c_t(i,j)$, are  dictated by the topology of the communication network,  with non-zero values only if  two agents $i$ and $j$ are connected at time $t$. They also need to satisfy the \emph{doubly stochastic} property in expectation, so that $\omega^i_{t}$ reaches a \emph{consensual} value for all $i\in\cN$ if it converges. 
Then, each agent $i$ updates its policy following stochastic policy gradient given by \eqref{eq:policy_gradient_thm} in the actor step, using its own Q-function estimate $Q_{{\theta}}(\cdot,\cdot;\omega^i_t)$.  
A variant algorithm  is also introduced in  \cite{zhang2018fully}, relying on not the Q-function, but the state-value function approximation, to estimate the global advantage function. 

With these in mind, almost sure convergence  is established    in \cite{zhang2018fully} for these decentralized actor-critic algorithms, when     linear functions are  used for value function approximation.  
Similar ideas are also extended to the setting with continuous spaces \cite{zhang18cdc}, where deterministic policy gradient (DPG) method  is usually used. Off-policy exploration, namely a stochastic behavior policy, is required  for DPG, as the deterministic on-policy may not be explorative enough. However, in the multi-agent setting, as the  policies of other agents are unknown, the common off-policy approach for DPG \cite[\S 4.2]{silver2014deterministic} does not apply. Inspired by the expected policy gradient (EPG) method \cite{ciosek2018expected} which unifies stochastic PG (SPG) and DPG, we develop an algorithm that remains on-policy, but reduces the variance of general SPG  \cite{zhang18cdc}. In particular, we derive the multi-agent version of EPG, based on which we  develop  the actor step that can be implemented in a decentralized fashion, while the critic step still follows  \eqref{equ:MARL_critic_1}. Convergence of the algorithm is then also established when linear function approximation is used \cite{zhang18cdc}. In the same vein, \cite{suttle2019multi} considers the extension of \cite{zhang2018fully} to an off-policy  setting, building upon the  emphatic temporal differences (ETD) method for the critic \cite{sutton2016emphatic}. By incorporating the analysis of ETD($\lambda$) \cite{yu2015convergence} into \cite{zhang2018fully}, almost sure convergence guarantee has also been established. Another off-policy algorithm for the same setting is proposed concurrently by  \cite{zhang2019distributed}, where agents do not share their estimates of value function. Instead, the agents aim to reach consensus over the global optimal policy estimation. Provable convergence is then established for the algorithm, with  a  local critic and a consensus actor.   


RL for decentralized networked agents has also been investigated in  \emph{multi-task}, in addition to the multi-agent, settings.  In some sense, the former can be regarded as a simplified version of the latter, where each agent deals with an \emph{independent MDP} that is not affected by other agents, while the goal is still to learn the  optimal joint policy that accounts for the average reward of all agents. 
\cite{pennesi2010distributed} proposes a distributed actor-critic algorithm,  assuming that the states, actions, and rewards are all local to each agent. Each agent performs a local TD-based critic step, followed by a consensus-based actor step that follows the gradient calculated using  information exchanged from the neighbors. Gradient of the average return is then proved to converge to zero as the iteration goes to infinity. \cite{macua2017diff} has  developed  \emph{Diff-DAC}, another distributed actor-critic algorithm for this setting, from duality theory. The updates resemble those  in \cite{zhang2018fully}, but provide additional insights that   actor-critic is  actually an instance of
the dual ascent method for solving  a linear program. 

Note that all the aforementioned convergence guarantees are \emph{asymptotic}, i.e., the algorithms converge as the iteration  numbers go to infinity, and are restricted to the case with linear function approximations. This fails to quantify the performance when  finite iterations and/or samples are used, not to mention when nonlinear functions such as deep neural networks are utilized. As an initial step toward \emph{finite-sample analyses} in this setting with  more \emph{general} function approximation, we consider   in \cite{zhang2018finite} the \emph{batch} RL algorithms \cite{lange2012batch}, specifically, decentralized variants of the fitted-Q iteration (FQI) \cite{riedmiller2005neural,antos2008fitted}. Note that we focus on FQI since it  motivates the celebrated deep Q-learning algorithm \cite{mnih2015human} when  deep neural networks are used for function approximation. 
We study FQI variants for both the cooperative setting with networked agents, and the competitive setting with two teams of such networked agents (see \S\ref{sec:value_comp} for more details). In the former setting,   all agents cooperate to iteratively update the global Q-function estimate, by fitting nonlinear least squares with the target values as the responses. In particular, let $\cF$ be the  function class for Q-function approximation, $\{(s_j,\{a^i_j\}_{i\in\cN},s_{j}')\}_{j\in[n]}$ be the batch transitions dataset of size $n$  available to all agents, $\{r^i_j\}_{j\in[n]}$ be the \emph{local} reward samples private to each agent, and $y^i_j=r^i_j+\gamma\cdot \max_{a\in\cA}Q^i_t(s_{j}',a)$  be the local target value, where $Q^i_t$ is agent $i$'s Q-function estimate at iteration  $t$. Then,  all agents cooperate to find a common Q-function estimate by  solving
\#\label{equ:FQI_coop}
\min_{ f\in \cF}~~ \frac{1}{N}\sum_{i\in\cN}\frac{1}{2n} \sum_{j=1}^n \bigl [ y^i_j - f(s_j, a_j^1, \cdots,a_j^N) \bigr ]^2.  
\#
Since $y^i_j$ is only known to agent $i$, the problem in \eqref{equ:FQI_coop} aligns with the formulation of \emph{distributed/consensus optimization} in \cite{nedic2009distributed,agarwal2011distributed,jakovetic2011cooperative,tu2012diffusion,hong2017stochastic,nedic2017achieving}, whose global optimal  solution can be achieved by the algorithms therein, if  $\cF$ makes  $\sum_{j=1}^n  [ y^i_j - f(s_j, a_j^1, \cdots,a_j^N)  ]^2$ convex for each $i$. This is indeed the case if $\cF$ is a linear function class. Nevertheless, with only a finite iteration of  distributed optimization algorithms (common in practice), agents may not reach exact consensus, leading to an error of each agent's Q-function estimate away from the actual optimum of \eqref{equ:FQI_coop}. Such an error also exists when nonlinear function approximation is used.  Considering this error caused by decentralized computation, we follow the \emph{error propagation} analysis stemming  from single-agent batch RL \cite{munos2007performance,munos2008finite,antos2008fitted,antos2008learning,farahmand2010error}, to establish the finite-sample performance of the proposed algorithms, i.e., how the accuracy of the algorithms output   depends on the function class $\cF$, the number of samples within each iteration $n$, and the number of iterations $t$. 


%The independent MDP case \cite{macua2017diff}, \cite{pennesi2010distributed} (where all state and action spaces are local, and the  each agent's action does not change the other agents'  transition models.)

\vspace{6pt}
\noindent{\bf Policy Evaluation}
\vspace{3pt}

Aside from control, a series of algorithms  in this setting focuses  on the policy evaluation task only, namely, the critic step of the actor-critic algorithms. With the policy fixed,  this task embraces a neater formulation,  as the sampling  distribution becomes stationary, and the objective becomes convex under linear function approximation. This facilitates the  finite-time/sample analyses, in contrast to  most control algorithms with only  {asymptotic guarantees}.   Specifically, under joint policy $\pi$, suppose each agent parameterizes the value function by a linear function class $\{V_\omega(s):=\phi^\top(s)\omega:\omega\in\RR^d\}$, where $\phi(s)\in\RR^d$ is the  feature vector at $s\in\cS$, and $\omega\in\RR^d$ is the vector of parameters. For notational convenience, let $\Phi:=(\cdots;\phi^\top(s);\cdots)\in\RR^{|\cS|\times d}$,  $D=\diag[\{\eta_\pi(s)\}_{s\in\cS}]\in\RR^{|\cS|\times|\cS|}$ be a diagonal matrix constructed using the   state-occupancy measure $\eta_\pi$, $\bar R^\pi(s)=N^{-1}\cdot \sum_{i\in\cN}R^{i,\pi}(s)$, where $R^{i,\pi}(s)=\EE_{a\sim\pi(\cdot\given s),s'\sim P(\cdot\given s,a)}[R^i(s,a,s')]$, and $P^\pi\in\RR^{|\cS|\times|\cS|}$ with the $(s,s')$ element being $[P^\pi]_{s,s'}=\sum_{a\in\cA}\pi(a\given s)P(s'\given s,a)$. 
Then, the objective of all agents is to jointly minimize the mean square projected Bellman error (MSPBE) associated with the team-average reward, i.e., 
\#\label{equ:MSPBE_joint}
\min_{\omega}\quad\texttt{MSPBE}(\omega):=\big\|\Pi_{\Phi}\big(V_\omega-\gamma P^\pi V_\omega -\bar R^\pi\big)\big\|^2_D=\big\|A\omega-b\big\|^2_{C^{-1}},
\#
where $\Pi_{\Phi}:\RR^{|\cS|}\to \RR^{|\cS|}$ defined as $\Pi_{\Phi}:=\Phi(\Phi^\top D\Phi)^{-1}\Phi^\top D$ is the projection operator onto subspace $\{\Phi\omega:\omega\in\RR^d\}$, $A:=\EE\{\phi(s)[\phi(s)-\gamma \phi(s')]^\top\}$, $C:=\EE[\phi(s)\phi^\top(s)]$, and $b:=\EE[\bar R^\pi(s)\phi(s)]$. By replacing the expectation with samples and using the 
Fenchel duality, the finite-sum  version of \eqref{equ:MSPBE_joint} can be re-formulated as a distributed saddle-point problem 
\#\label{equ:MSPBE_joint_finite_sum}
\min_{\omega}\max_{\lambda^i,i\in\cN}\quad\frac{1}{Nn}\sum_{i\in\cN}\sum_{j=1}^n 2(\lambda^i)^\top A_j\omega -2(b^i_j)^\top \lambda^i-(\lambda^i)^\top C_j\lambda^i
%\texttt{MSPBE}(\omega):=\big\|\Pi_{\Phi}\big(V_\omega-\gamma P^\pi V_\omega -\bar R^\pi\big)\big\|^2_D
\#
where $n$ is the data size, $A_j,C_j$ and $b^i_j$ are empirical estimates of $A,C$ and $b^i:=\EE[R^{i,\pi}(s)\phi(s)]$, respectively, using sample $j$.  
Note that \eqref{equ:MSPBE_joint_finite_sum} is convex in $\omega$ and concave in $\{\lambda^i\}_{i\in\cN}$. The use of   MSPBE as an objective is standard in multi-agent policy evaluation \cite{macua2015distributed,lee2018primal,wai2018multi,doan2019finitea,doan2019finiteb}, and the idea of  saddle-point reformulation  has been adopted in \cite{macua2015distributed,lee2018primal,wai2018multi,cassano2018multi}. Note that in \cite{cassano2018multi}, a variant of MSPBE, named H-truncated $\lambda$-weighted MSPBE, is advocated, in order to control the bias of the solution deviated from the actual mean square Bellman error minimizer.  

With the formulation \eqref{equ:MSPBE_joint} in mind, \cite{lee2018primal} develops a distributed variant of the gradient TD-based method, with asymptotic   convergence established using the ordinary differential equation (ODE) method.  In \cite{wai2018multi},  a double averaging scheme that combines the dynamic consensus  \cite{qu2017harnessing} and the SAG algorithm \cite{schmidt2017minimizing} has been proposed to solve the saddle-point problem \eqref{equ:MSPBE_joint_finite_sum} with a linear rate. \cite{cassano2018multi} incorporates the idea of   variance-reduction, specifically, AVRG in \cite{ying2018convergence}, into  gradient TD-based policy evaluation. Achieving the same linear rate as \cite{wai2018multi}, three  advantages are claimed in  \cite{cassano2018multi}: i)  data-independent memory requirement; ii) use of eligibility traces \cite{singh1996reinforcement}; iii) no need for  synchronization  in sampling. More recently, standard TD learning \cite{tesauro1995temporal}, instead of gradient-TD, has been generalized to this MARL setting, with special focuses on  finite-sample analyses \cite{doan2019finitea,doan2019finiteb}.  
Distributed TD($0$) is first studied in \cite{doan2019finitea}, using the proof techniques originated in \cite{pmlr-v75-bhandari18a}, which requires a projection on the iterates, and the data samples to be independent and identically distributed (i.i.d.). Furthermore, motivated by the recent progress in \cite{pmlr-v99-srikant19a}, finite-time performance of the  more general  distributed TD($\lambda$) algorithm is provided in \cite{doan2019finiteb}, with neither projection nor i.i.d. noise assumption needed.  



Policy evaluation for networked agents has also been investigated under the setting of \emph{independent} agents interacting with \emph{independent} MDPs.   \cite{macua2015distributed} studies off-policy evaluation based on  the importance sampling technique. With no coupling among MDPs, an agent  does not need to know the actions of the other agents.  Diffusion-based distributed GTD is then proposed, and is shown to convergence in the mean-square sense with a sublinear rate. In  \cite{stankovic2016multi}, two variants of the TD-learning, namely, GTD2 and TDC \cite{sutton2009fast}, have been  designed for this setting, with weak convergence proved by the general stochastic approximation theory in \cite{stankovic2016distributed}, when agents are connected by a time-varying communication network. Note that \cite{cassano2018multi} also considers the independent MDP setting, with the same results established as the actual MARL setting. 


%Each agent maintains a local copy of $\theta$, 
% \issue{2019.09.22}
%  
% serves as the common cost  function to minimize
%
%
%\issue{XXXX} easier and thus hope to get better results. 


%introduce the algorithms first, including both for \emph{control} and for \emph{policy evaluation} \cite{wai2018multi,cassano2018multi,doan2019finitea,doan2019finiteb},.

%Also, note that several distributed results (over network) assume ``independent MDP'', for policy evaluation \cite{macua2015distributed,stankovic2016distributed,stankovic2016multi}. 

\vspace{6pt}
\noindent{\bf Other Learning Goals}
\vspace{3pt}

Several other learning goals have also been explored for decentralized MARL with networked agents. \cite{zhang2016data} has considered  the \emph{optimal consensus}  problem, where each agent over the network tracks the states of its neighbors' as well as a leader's, so that the  consensus error is minimized by the joint policy. A policy iteration algorithm is then introduced,  followed by a practical  actor-critic algorithm using neural networks for function approximation. A similar consensus error objective is also adopted in \cite{zhang2018model}, under  the name of \emph{cooperative multi-agent graphical games}. A centralized-critic-decentralized-actor scheme is utilized for developing off-policy RL algorithms. 

Communication efficiency, as a key ingredient in the algorithm design for this setting, has drawn increasing attention recently \cite{chen2018communication,ren2019communication,lin2019communication}.  Specifically, \cite{chen2018communication} develops Lazily
Aggregated Policy Gradient (LAPG), a distributed PG algorithm that can reduce  the communication rounds between the agents and a central controller, by judiciously designing communication trigger rules. \cite{ren2019communication} addresses the same policy evaluation problem as \cite{wai2018multi}, and develops a   hierarchical distributed algorithm by proposing a mixing matrix different from the doubly stochastic one used in \cite{zhang2018fully,wai2018multi,lee2018primal}, which allows unidirectional information exchange among agents to save  communication. In contrast, the distributed actor-critic algorithm in  \cite{lin2019communication} reduces the communication by transmitting only   one scalar entry of its state vector at each iteration, while preserving provable  convergence as in \cite{zhang2018fully}. 



%mention the other related research, e.g., comm. efficient., under attack/with robustness, 


%\begin{itemize}
%%	\item Independent Q learning works in this case, with convergence.
%%\item Learning in Markov potential games \cite{valcarcel2018learning}, survey or modeling of MPG \cite{gonzalez2013discrete,zazo2016dynamic} 
%\item MARL with Networked agents: Q-learning \cite{kar2013cal};  Actor-Critic \cite{zhang2018fully}, and all follow-ups, including \emph{asymptotic} \cite{lee2018primal}  and \emph{finite-time} analyses for policy evaluation \cite{wai2018multi,cassano2018multi,doan2019finitea,doan2019finiteb}, and actor-critic \cite{suttle2019multi} and consensus on actor   \cite{zhang2019distributed}. 
%%\item RNN / LSTM-based method (DRQN)
%%\item Learning to communicate and its followup work
%\item Cooperative network/graphical games: \cite{zhang2018model} 
%\item Communication-efficient MARL \cite{chen2018communication,ren2019communication,lin2019communication}
%%\item Mean-field control / swarm system (a large number of homogeneous agents), or mean-field teams.
%%\item Actor-critic fictitious play for \emph{multi-stage} MARL \cite{perolat2018actor}. 
%%\item Fictitious play also works for cooperation, namely, is shown to converge for cooperative, potential (normal-form) games \cite{hofbauer2002global}.
%\end{itemize}



\subsubsection{Partially Observed Model}\label{subsec:partially_observed}

We complete the overview for cooperative settings  by briefly introducing a class of significant but challenging models where agents are faced with partial  observability. Though common in practice, theoretical analysis of MARL algorithms in this setting is still relatively scarce, in contrast to the aforementioned fully observed settings.  
In general, this setting can be modeled by a decentralized POMDP (Dec-POMDP) \cite{oliehoek2016concise}, which shares almost all elements such as the reward function and the transition model, as the MMDP/Markov team model in \S\ref{subsec:Markov_Games}, except that each agent now only has its local observations of the system state $s$. With no accessibility  to other agents' observations, an  individual agent cannot maintain a global belief state, the sufficient statistic for decision making in single-agent POMDPs. Hence, Dec-POMDPs have been known to be  NEXP-complete \cite{bernstein2002complexity},    requiring super-exponential time to solve in the worst case.  

There is an increasing interest in developing planning/learning algorithms for Dec-POMDPs. 
Most of the algorithms are based on the \emph{centralized-learning-decentralized-execution} scheme. In particular, the  decentralized problem is first reformulated as a centralized one, which can be solved at a central controller with (a simulator that generates) the observation data of all agents. The policies are then optimized/learned using data, and distributed to all agents for  execution.   
Finite-state controllers (FSCs) are commonly used to represent  the local policy at each agent \cite{bernstein2009policy,amato2010optimizing}, which map local observation histories to actions. 
A Bayesian nonparametric approach is proposed in \cite{liu2015stick} to   determine the controller size of variable-size FSCs. 
To efficiently solve the centralized problem, a series of \emph{top-down}  algorithms have been proposed. In \cite{oliehoek2014dec},   the Dec-POMDP is converted  to  \emph{non-observable MDP} (NOMDP), 
a kind of centralized sequential decision-making problem, which is then  addressed by some  heuristic  tree search algorithms. 
As an extension of the NOMDP conversion, 
\cite{dibangoye2016optimally,dibangoye2018learning} convert Dec-POMDPs to   \emph{occupancy-state} MDPs (oMDPs), where the occupancy-states are distributions over hidden states and joint histories of observations. As the value functions of  oMDPs enjoy the piece-wise linearity and convexity, both tractable planning \cite{dibangoye2016optimally} and value-based learning \cite{dibangoye2018learning} algorithms have been developed. 
  
To further improve  computational  efficiency, 
several sampling-based planning/learning algorithms have also been proposed.  {In particular, } Monte-Carlo sampling {with} policy iteration and {the expectation-maximization} algorithm, are proposed in \cite{wu2010rollout} and \cite{wu2013monte},  respectively. 
Furthermore, Monte-Carlo tree search   has been applied {to} special classes of Dec-POMDPs, such as  multi-agent POMDPs \cite{amato2015scalable} and multi-robot active perception \cite{best2018dec}.  
In addition, policy gradient-based algorithms can also be developed for this centralized learning scheme \cite{dibangoye2018learning}, with a centralized critic and a decentralized actor.  
Finite-sample analysis can also be   established under this scheme \cite{amato2009achieving,banerjee2012sample}, for tabular settings with finite state-action spaces. 

Several attempts have also been  made to  enable \emph{decentralized learning} in Dec-POMDPs. When the agents share some common information/observations, \cite{nayyar2013decentralized} proposes to reformulate the problem as a centralized POMDP, with the common information being the observations of a \emph{virtual} central controller. This way, the centralized POMDP can be solved individually by each agent. In \cite{arabneydi2015reinforcement}, the reformulated POMDP has been  approximated by finite-state MDPs with exponentially decreasing approximation error, which are then solved by  Q-learning. Very  recently, \cite{zhang2019online} has developed a tree-search based algorithm for solving this centralized POMDP, which, interestingly, echoes back the heuristics for solving Dec-POMDPs directly as in \cite{amato2015scalable,best2018dec}, but with a more solid theoretical footing. Note that in both \cite{arabneydi2015reinforcement,zhang2019online}, a common random number generator  is used for all agents, in order to avoid communication among agents and enable a decentralized learning scheme.   
 
 

%\issue{Mention the difference of the model first, then review alg. from ACC.}

% Dec-POMDP solvers (briefly) with theory:   MARL with partial information sharing \cite{arabneydi2015reinforcement,zhang2019online} for Dec-POMDP, from a decentralized control perspective.  Introduce MCTS here. MCTS contains a few components. Action selection is usually done using UCB1 algorithm. Theoretical results: Cite Qiaomin (nonasymptotic)  and David Silver (asymptotic) \issue{just mention a little bit, since we do not have models for Dec-POMDP.}. ,  Oliehoak's works with theories, etc; 
 
 
 
\subsection{Competitive Setting}

\iffalse
%\begin{comment}
\noindent{XXXXXXXXXXXXX}
\issue{@Zhuoran: Plz check out and clean the comments you left.}
XXXXXXXXXXXXXX

Maybe also survey online learning methods / repeated zero-sum matrix games. There is a large literature. 

Also related to robust MDP, risk sensitive MDP, and constrained MDP. 

{\color{red} What I haven't covered: 

\begin{itemize}
\item General-sum: static-mixed setting in Figure 1 in \cite{bu2008comprehensive}, but they are all just small-scale. ``Most' theories are for static setting only, while for dynamics settings, most are heurestics.''
\item General-sum: Mention \cite{bowling2002multiagent}, which compares the standard minimax-Q and Nash-Q, which are both convergent, but not rational. Instead, it proposes IGA, WoLF-IGA, which are gradient based methods, with convergence to two-player, two-agent, general-sum repeated matrix games.  
Then summarize the all these have theory for tabular cases only. This WoLF idea is applied to actual ``dynamic games'' by WoLF-PHC also in  \cite{bowling2002multiagent}, without theoretical guarantees.
\item LQ case,  Q-learning  \cite{al2007model},  heuristic DP and dual-heuristic DP \cite{al2007adaptive}
%\item Model-based: E3 for single controller SG (SCSG) (Brafman and Tennenholtz, 2000);
%R-Max for general zero-sum SG (Brafman and Tennenholtz, 2002)
\item Policy gradient for zero-sum continuous games/saddle-point problems  \cite{mazumdar2019finding,jin2019minmax,zhang2019policyb}
\item {Remember to cite related surveys}: Review of zero sum games \cite{bovsansky2016algorithms}
\end{itemize}


\issue{memory notes}
\begin{itemize}
	\item Other formalisms/frameworks of MARL: e.g., partially-observed setting for non-cooperative, say,  POSG; just normal-form games, say, fictitious-play with convergence \cite{hofbauer2002global}; 
repeated normal game setting (multi-agent learning)  \cite{claus1998dynamics,bowling2002multiagent,conitzer2007awesome,banerjee2003adaptive,singh2000nash,bowling2005convergence,chakraborty2014multiagent}, seminal works Sec. $3.1$ of \cite{hernandez2018multiagent}, [136-146], and recent progress \cite{song2019convergence} and references therein; simultaneous move games by MCTS \cite{lisy2013convergence}. 
\end{itemize}

}

\noindent{XXXXXXXXXXXXXXXXXXXXXXXXXXXXXXXXXXXXXXXXXXXXXXXXXXXXXXXXXXXXXXXXX}
%\end{comment}
 
 \fi 
 
 
 
 Competitive settings are usually modeled as \emph{zero-sum} games. Computationally, 
  there exists a great barrier between solving  two-player and multi-player zero-sum games. 
 In particular, even the  simplest   three-player   matrix games, are known to be 
 PPAD-complete \cite{papadimitriou1992inefficient,daskalakis2009complexity}.
 Thus, most existing results on competitive MARL  focus on two-player zero-sum games, with $\cN = \{1, 2\}$ and $R^1 + R^2 = 0$ in Definitions \ref{def:Markov_Game}  and \ref{def:ext_form_game}. 
 In the rest of this section, we review methods that provably find  a Nash (equivalently, saddle-point)  equilibrium in two-player Markov   or  extensive-form games.  
 The existing algorithms can mainly  be categorized into two classes: value-based  and policy-based approaches, 
 which are introduced separately in the sequel.
 
  
 \subsubsection{Value-Based Methods} \label{sec:value_comp}
 Similar as in single-agent MDPs, value-based methods aim to find an optimal value function  from which the joint  Nash equilibrium policy can be extracted.  
 Moreover, the optimal value function is known to be the unique fixed point of a Bellman operator, which can be obtained via dynamic programming type methods.
  
 Specifically, for simultaneous-move Markov games,  the value function defined in \eqref{eq:V_pi} satisfies  $V_{\pi^1, \pi^2}^1 = - V^2 _{\pi^1, \pi^2}   $. 
 and thus any   Nash equilibrium  $\pi^* = (\pi^{1,*}, \pi^{2,*} )$ satisfies 
 \#\label{eq:zero-sum-eq}
 V  _ {\pi^1, \pi^{1,*}} ^1 (s) \leq V  _ {\pi^{1,*}, \pi^{2,*}} ^1 (s)  \leq V^1   _ {\pi^{1,*}, \pi ^2} (s), \qquad \text{for any~}\pi = (\pi^1, \pi^2) ~\text{and}~s\in \cS. 
 \#
Similar as the  minimax theorem  in normal-form/matrix zero-sum games \cite{von2007theory}, for two-player zero-sum Markov games with finite state and action spaces,  one can define the  optimal value function $V^* \colon \cS \rightarrow \RR$  as
\#\label{eq:minimax}
V^* =   \max_{\pi^1} \min _{\pi^2} V^1 _{\pi^1, \pi^2}  = \min _{\pi^2} \max_{\pi^1}  V^1 _{\pi^1, \pi^2},
\#
Then \eqref{eq:zero-sum-eq} implies that 
$V_{\pi^{1,*}, \pi^{2,*}} ^1$  coincides with $V^*$ and any pair of policies $\pi^1$ and $\pi^2$ that attains  
 the supremum and infimum in \eqref{eq:minimax} constitutes a Nash equilibrium.
 Moreover, similar to MDPs, \cite{shapley1953stochastic} shows that $V^*$ is the unique solution of a Bellman equation and a Nash equilibrium can be constructed based on $V^*$. 
 Specifically,
 for any $V \colon \cS \rightarrow \RR$ and any $s \in \cS$, we define 
 \#\label{eq:Q_V}
 Q_V(s,a^1, a^2)  = \EE_{s' \sim \cP(\cdot \given s, a^1, a^2)} \big [R^1 (s, a^1, a^2, s' ) + \gamma \cdot V(s') \bigr ], 
 \#
where $Q_V(s,\cdot,\cdot)$ can be regarded as a matrix in $\RR^{| \cA^1|  \times | \cA^2|}$. Then we define the Bellman operator  $\cT^*$ by solving a  matrix zero-sum game regarding  $Q_V(s, \cdot, \cdot )$ as the payoff matrix, i.e., for any $s\in \cS$,  one can define 
 \#\label{eq:bellman_eq}
( \cT^* V) (s) = \texttt{Value} \bigl[ Q_V(s,\cdot , \cdot ) \bigr ]  = \max_{u \in \Delta(\cA^1) }\min _{v\in \Delta(\cA^2) }  \sum_{a\in \cA^1} \sum_{b\in \cA^2} u_{a} \cdot  v_b  \cdot Q _V(s,a, b), 
 \# 
 where we use $ \texttt{Value} (\cdot) $ to denote the optimal value of a matrix zero-sum game, which can be obtained by solving a  linear program  \cite{vanderbei2015linear}. 
 Thus,  the Bellman operator $\cT^*$ is $\gamma$-contractive in the $\ell_{\infty}$-norm and $V^*$ in \eqref{eq:minimax} is the unique solution to the Bellman equation $V = \cT^* V$. Moreover, letting $ p_1(V), p_2(V)$ be any solution to the 
  optimization problem in  \eqref{eq:bellman_eq}, we have that 
 $ \pi^* = (p_1(V^*), p_2(V^*) ) $ is a Nash equilibrium specified by Definition \ref{def:NE}. Thus, 
based on the Bellman operator $\cT^*$, \cite{shapley1953stochastic} proposes  the value iteration algorithm, which creates a sequence of value functions $\{ V_{ t }\}_{t \geq 1} $ satisfying $V_{t+1} = \cT ^* V_{t} $ that converges to $V^*$ with a linear rate. Specifically, we have 
 \$
 \| V_{t+1}  - V^* \|_{\infty} = \| \cT^* V_{ t  } -  \cT^* V^*  \|_{\infty} \leq \gamma \cdot  \|V_{t} -    V^* \|_{\infty} \leq \gamma ^{t+1} \cdot  \| V_{0} - V^* \|_{\infty}  . 
 \$
 In addition, a value iteration update  can be decomposed into the two steps.  In particular, 
  letting $\pi^1  $ be any   policy of player $1$ and  $V$ be any value function, we define Bellman operator $\cT^{\pi^1 }$ by 
 \#\label{eq:bellman_oper2}
 ( \cT^{\pi^1} V)(s) =  \min _{v\in \Delta(\cA^2) }  \sum_{a\in \cA^1} \sum_{b\in \cA^2} \pi^1(a\given s ) \cdot v_b\cdot  Q _V(s,a, b),
 \# 
 where $Q_V$ is defined in \eqref{eq:Q_V}. 
 Then we can equivalently  write a  value iteration update as 
\#\label{eq:value_iteration}
\mu_{t+1} = p_1(V_{t}) \qquad \text{and}\qquad   V_{ t+1 }= \cT^{\mu_{ t+1}} V_{t}. 
 \# 
 Such a decomposition 
 motivates the policy iteration  algorithm for two-player zero-sum games, which has been  studied in, e.g.,   \cite{hoffman1966nonterminating, van1978discounted, rao1973algorithms, patek1997stochastic, hansen2013strategy}, for different variants of such Markov games. 
 In particular,  from the perspective of player $1$, policy iteration creates a sequence $\{ \mu_{ t} , V_{ t }\}_{t\geq 0} $ satisfying  
 \#\label{eq:policy_iteration}
 \mu_{ t+1 } = p_1 ( V_{t  } ) \qquad \text{and} \qquad V_{ t+1 }  =(\cT^{\mu_{ t+1 }})^{\infty} V_{ t} ,
 \#
i.e.,    $V_{ t+1 } $ is the fixed point of   $\cT^{\mu_{ t+1 }}$. The updates for  player $2$ can be similarly constructed. By the definition in \eqref{eq:bellman_oper2}, the Bellman operator  $ \cT^{ \mu_{ t+1 }}$  is $\gamma$-contractive and its fixed point corresponds to  the value function associated with $(  \mu_{ t+1 }, \texttt{Br}  ( \mu_{ t+1 }) )$, where $ \texttt{Br}  ( \mu_{ t+1 }) $ is the best response policy of player $2$ when player $1$ adopts $ \mu_{ t+1 }$.  Hence, in each step of policy iteration, the player first finds an improved policy  $ \mu_{ t+1 }$ based on the current  function $V_{t}$, and then obtains a conservative value function by assuming that the opponent plays the best counter policy $ \texttt{Br} (\mu_{t+1}) $. It has been  shown in  \cite{hansen2013strategy} that the  value function sequence  $\{ V_{t} \}_{t\geq 0}$ monotonically increases to $V^*$  with a linear rate of convergence for turn-based zero-sum Markov games. 
 
 Notice  that both the  value and  policy iteration algorithms are model-based due to the need of computing the Bellman operator $\cT^{ \mu_{ t+1 }}$ in \eqref{eq:value_iteration} and \eqref{eq:policy_iteration}. 
By estimating the Bellman operator
via data-driven approximation,  
 \cite{littman1994markov} has proposed   minimax-Q learning, which extends the well-known  Q-learning algorithm  \cite{watkins1992q} for MDPs to zero-sum Markov games. In particular,  minimax-Q learning is an online, off-policy, and tabular method 
 which updates the         action-value function    $Q \colon \cS\times \cA \rightarrow \RR$ based on transition data $\{(s_t, a_t, r_t, s_t')\}_{t\geq 0}$, where $s_t' $ is the next state following $(s_t, a_t)$ and $r_t$ is the reward.   
 In the $t$-th iteration, it only updates  the value  of $Q(s_t, a_t) $    and keeps other entries of $Q$   unchanged. Specifically, we have 
   \#\label{eq:minimaxQ}
   Q  (s_t,a^1_t, a^2_t  ) \leftarrow   ( 1- \alpha_t )  \cdot Q (s_t,a^1_t, a^2_t   ) +\alpha_t  \cdot \bigl \{  r_t + \gamma \cdot \texttt{Value} \big[ Q  (s_t', \cdot,  \cdot) \bigr ]  \bigr \} ,
   \#
   where $\alpha_t \in (0,1)$ is the stepsize.
 {As shown in \cite{szepesvari1999unified}, under  conditions similar to those for   single-agent Q-learning \cite{watkins1992q}, function $Q$ generated by \eqref{eq:minimaxQ}   
   converges to the optimal action-value function $Q^* = Q_{V^*}$ defined by combining \eqref{eq:minimax} and \eqref{eq:Q_V}.} 
   Moreover, with a slight abuse of notation, if we define the  Bellman operator $\cT^*$ for action-value functions by  
    \#\label{eq:minimax_Q}
   ( \cT^* Q )(s,a^1, a^2 )  = \EE _{  s' \sim \cP(\cdot \given s, a^1, a^2)} \big  \{ R^1 (s, a^1, a^2, s' ) + \gamma \cdot {\tt{Value}}\bigl[  Q(s', \cdot, \cdot) \bigr ] \bigr \}, 
   \# 
  then we have  
   $Q^*$ 
  as the unique fixed point of $\cT^*$. 
  Since the target value 
  $ r_t + \gamma \cdot   \texttt{Value} [ Q  (s_t', \cdot,  \cdot) ] $ in \eqref{eq:minimaxQ} is an  estimator of $(\cT^* Q) (s_t, a_t^1, a_t^2)$, 
   minimax-Q learning can be viewed as a stochastic approximation algorithm for computing the fixed point of $\cT^*$. 
   Following \cite{littman1994markov}, 
  minimax-Q learning has been further extended to the function approximation setting  where $Q$ in \eqref{eq:minimaxQ} is approximated by a class of parametrized functions.  However, convergence guarantees for this minimax-Q learning with even linear function approximation have not been well understood. 
%  In particular,  
%   \cite{lagoudakis2002value,zou2019finite} establish the convergence of minimax-Q learning with linear function approximation and temporal-difference updates \cite{sutton1987temporal}.    
   Such a linear value function approximation also applies to a significant class  of zero-sum MG instances with continuous state-action spaces, i.e., linear quadratic (LQ) zero-sum games \cite{bacsar1995h,al2007adaptive,al2007model}, where  the reward function is quadratic with respect to the states and actions, while the transition model follows linear dynamics. In this setting, Q-learning based algorithm can be  guaranteed  to converge  to the NE \cite{al2007model}. 
   
   
   
   To embrace general function classes, the framework of batch RL \cite{lagoudakis2002value,munos2007performance,munos2008finite,antos2008fitted,antos2008learning,farahmand2016regularized} can be adapted to  the multi-agent settings, as in the recent works  \cite{yang2019theoretical,zhang2018finite}.      As mentioned  in \S\ref{subsubsec:networked_paradigm} for cooperative batch MARL,   each agent  iteratively  updates the Q-function by fitting least-squares using  the target values.  Specifically, let $\cF$ be the function class of interest and let $  \{ (s_i, a_i^1, a_i^2, r_i, s_i' ) \}_{i\in[n]}$ be the dataset. For any $t \geq 0$,   let $Q_{t}$ be the current iterate in the $t$-th iteration, and  define $y_i = r_i + \gamma \cdot \texttt{Value} [ Q_{t} (s_i', \cdot , \cdot )]$ for all $i \in [n]$. Then we update $ Q_{t}$   by solving a least-squares regression problem in $\cF$, that is, 
  \#\label{eq:FQI}
  Q_{t+1} \leftarrow \argmin_{ f\in \cF} \frac{1}{2n} \sum_{i=1}^n \bigl [ y_i - f(s_i, a_i^1, a_i^2 ) \bigr ]^2.
  \#
  In such a two-player zero-sum Markov game setting, 
a finite-sample error bound on the Q-function estimate  is established in    \cite{yang2019theoretical}. 


Regarding other finite-sample analyses, 
very recently, \cite{jia2019feature} has studied   zero-sum turn-based stochastic games (TBSG), a simplified zero-sum MG when the transition  model is embedded in some feature space and a generative model is available. Two Q-learning based  algorithms have been  proposed and analyzed  for this setting.  \cite{sidford2019solving} has proposed   algorithms that achieve near-optimal sample complexity for general zero-sum TBSGs with a generative model, by extending the previous near-optimal Q-learning algorithm for MDPs \cite{sidford2018near}.   
In the online setting, where the learner controls only one of the players that plays against an arbitrary opponent,  
\cite{wei2017online} has proposed UCSG, an algorithm for the  \emph{average-reward} zero-sum  MGs, using the principle of \textit{optimism in the  face of uncertainty} \cite{auer2007logarithmic, jaksch2010near}.  UCSG is shown to achieve a sublinear regret  compared to the game value when competing with an arbitrary opponent, and also achieve $\tilde O(\text{poly}(1/\epsilon))$ sample complexity if the opponent plays an optimistic best response. 


 Furthermore, when it comes to zero-sum games  with imperfect information, \cite{kollery1994fast, von1996efficient, koller1996efficient, von2002computing} have proposed to transform  extensive-form games into normal-form games using the 
 \textit{sequence form} representation, which enables equilibrium  finding via linear programming. 
 In addition,   by lifting the state space to  the space of belief states,  \cite{parr1995approximating,rodriguez2000reinforcement,hauskrecht2000value, hansen2004dynamic,buter2012dynamic} have applied dynamic programming methods to zero-sum stochastic games. 
 Both of these approaches  guarantee  finding of a Nash equilibrium but are only efficient for small-scale problems. 
 Finally,  MCTS with UCB-type action selection rule  \cite{chang2005adaptive,kocsis2006bandit,coulom2006efficient} can also be applied to two-player turn-based games with incomplete  information \cite{kocsis2006bandit, cowling2012information, teraoka2014efficient,whitehouse2014monte,  kaufmann2017monte}, which lays the foundation for the recent success of deep RL for the game of Go \cite{silver2016mastering}. Moreover, these methods are 
 shown to converge to the minimax solution of the game, thus can be viewed as a counterpart of minimax-Q learning with Monte-Carlo sampling. 
 %and require  the knowledge of the model. 

  
 

    
  \subsubsection{Policy-Based Methods} \label{sec:policy_comp}
 
 Policy-based reinforcement learning methods introduced in \S\ref{sec:policyRL}
 can  also be extended to the multi-agent setting. Instead of finding the fixed point of the Bellman operator, a fair amount of methods only focus on a single agent and aim to  maximize the expected return  of that agent, disregarding  the other agents' policies.
 Specifically, 
 from the perspective of a single agent, the environment is time-varying as the other agents also  adjust their policies. Policy based methods aim to achieve the  optimal performance when other agents play arbitrarily by minimizing the 
   \textit{(external) regret}, that is, find a sequence of actions that perform nearly as well as  the  optimal fixed policy   in hindsight.
 An algorithm with negligible average overall    regret is  called \textit{no-regret} or \textit{Hannan consistent}  \cite{hannan1957approximation}.
 Any Hannan consistent algorithm is  known to have the following two desired properties in repeated normal-form games. First, when other agents adopt stationary policies, the  time-average policy constructed by the  algorithm converges to the best response policy (against the ones used by the other agents). Second, more interestingly, in two-player zero-sum games, when both players adopt Hannan consistent algorithms and both their  average overall regrets are no more than $\epsilon$,    their time-average policies constitute a $2\epsilon$-approximate Nash equilibrium \cite{blum2007learning}. 
  Thus, 
  any Hannan consistent single-agent reinforcement learning algorithm can be applied to find the Nash equilibria of  
   two-player zero-sum games via \emph{self-play}. Most of    these methods belong to one of  the following two families: 
  fictitious play \cite{brown1951iterative, robinson1951iterative}, and counterfactual regret minimization \cite{zinkevich2008regret}, which will be summarized below. 
 
 
  
  Fictitious play  is a classic algorithm studied in game 
theory, where  the players  play the game repeatedly and each player adopts a policy that  best responds to the average policy of the other agents. 
This method was originally proposed for solving 
 \textit{normal-form games}, which are a simplification of the Markov  games defined in Definition \ref{def:Markov_Game} with $\cS  $  being a singleton and $\gamma = 0$.  
 In particular, for  any joint policy $\pi \in \Delta(\cA)$    of the  $N$ agents, we let $\pi^{-i}$ be the marginal policy of all players except player $i$. For any $t\geq 1$, suppose the agents have played $\{ a_\tau \colon.1\leq  \tau   \leq t \} $ in the first $t$ stages. We define $x_{t }$ as the empirical distribution of  $\{ a_\tau \colon.1\leq  \tau   \leq t \} $, i.e., $x_t(a) = t^{-1} \sum_{\tau =1}^t\mathbbm{1} \{ a_t = a \} $ for any $a \in \cA$.  Then, in the $t$-th stage, each agent $i$ takes action $a_t^i \in \cA^i $ according to the best response policy against $x_t^{-i} $.  In other words, each agent plays the best counter policy against  the policy of the other agents inferred from  history data. 
 Here, for any $\epsilon >0$ and any $\pi \in \Delta(\cA)$, we denote by $\texttt{Br}_{\epsilon} (\pi^{-i})$ the $\epsilon$-best response policy of player $i$, which satisfies 
 \#\label{eq:br_policy}
 R^i \bigl  ( \texttt{Br}_{\epsilon} (\pi^{-i}), \pi^{-i} \bigr ) \geq \sup_{\mu\in \Delta(\cA^i)} R^i(\mu, \pi^{-i})  -\epsilon.
 \#
 Moreover, we define $\texttt{Br}_{\epsilon} (\pi)$ as the joint policy $( \texttt{Br}_{\epsilon} (\pi^{-1}) , \ldots,  \texttt{Br}_{\epsilon} (\pi^{-N}) ) \in \Delta (\cA) $ and suppress the subscript $\epsilon$ in $\texttt{Br}_{\epsilon}$ if $\epsilon = 0$.
By this notation, regarding each $a \in \cA$ as a vertex of $\Delta( \cA)$, we can equivalently write the fictitious process as 
 \#\label{eq:DFP}
 x_t - x_{t-1} = (1/ t)  \cdot (a_t - x_{t-1} ), \qquad \text{where}\qquad a_t \sim \texttt{Br}( x_{t-1}  )  . 
 \#
 As $t\to\infty$, the updates  in  \eqref{eq:DFP}  can be approximately characterized by  a differential inclusion \cite{benaim2005stochastic} 
 \#\label{eq:CFP}
 \frac{\ud x(t)  }{\ud t} \in  \texttt{Br}( x(t) ) - x(t),
 \#
 which is known as the continuous-time fictitious play. Although it is well known that  the discrete-time fictitious play  in \eqref{eq:DFP} is not Hannan consistent \cite{hart2001general, young1993evolution}, it is shown in 
 \cite{monderer1997belief, viossat2013no} that the continuous-time fictitious play  in  \eqref{eq:CFP}   is Hannan consistent. 
 Moreover, using tools from stochastic approximation \cite{kushner2003stochastic, hart2001general},  
 various modifications of  discrete-time  fictitious play based on techniques such as  smoothing or stochastic perturbations have been  shown to converge to  the  continuous-time fictitious play   and are thus  Hannan consistent
 \cite{fudenberg1995consistency, hofbauer2002global,  leslie2006generalised, benaim2013consistency, li2018sampled}. As a result,  applying these methods with self-play provably finds a Nash equilibrium of a two-player zero-sum normal form game. 
 
 
 Furthermore, fictitious play methods have also been extended to RL settings without the model knowledge.  Specifically,  using  sequence-form representation,
 %\cite{kollery1994fast, von1996efficient, koller1996efficient, von2002computing},  
 \cite{heinrich2015fictitious}  has proposed the first fictitious play algorithm  for extensive-form games which is realization-equivalent to the \textit{generalized weakened fictitious play} \cite{leslie2006generalised} for normal-form games. 
 The pivotal insight is that a convex combination of normal-form policies can be written as a weighted convex combination of behavioral policies using realization  probabilities. 
 Specifically,  recall that  the set of information states of agent $i$ was denoted by $\cS^i$.  
 When the game has perfect-recall, each $s^i \in \cS^i$ uniquely defines a sequence $\sigma_{s^i}$ of actions played by agent $i$ for reaching state $s^i$.  Then any behavioral policy $\pi^i$ of agent $i$ induces a \textit{realization probability} $\texttt{Rp}(\pi^i;  \cdot  )$  for each  sequence $\sigma$ of agent $i$, which is defined by $\texttt{Rp}(\pi^i;  \sigma ) = \prod_{(\sigma_{s'},  a) \sqsubseteq \sigma } \pi^i( a \given s') $, where the product is taken over all $s' \in \cS^i$ and $a\in \cA^i$ such that   $(\sigma_{s'} , a) $ is a subsequence of $\sigma$. 
 Using the notation of realization probability, for any two behavioral policies  $\pi$ and $\tilde \pi$ of agent $i$,  the sum 
 \#\label{eq:mixture}
   \frac{\lambda \cdot \texttt{Rp} (\pi, \sigma_{s^i})  \cdot \pi( \cdot \given s^i) }  { \lambda \cdot \texttt{Rp} (\pi, \sigma_{s^i})  + (1- \lambda) \cdot \texttt{Rp} ( \tilde \pi  , \sigma_{s^i})}+  \frac{(1- \lambda) \cdot \texttt{Rp} ( \tilde \pi , \sigma_{s^i}) \cdot  \tilde \pi (\cdot \given s^i )}  { \lambda \cdot \texttt{Rp} (\pi, \sigma_{s^i})  + (1- \lambda) \cdot \texttt{Rp} ( \tilde \pi , \sigma_{s^i})}  , \qquad \forall s^i \in \cS^i, 
 \#
is the mixture policy of  $\pi$ and $\tilde \pi$ with weights $\lambda \in (0,1)$ and $1-\lambda$, respectively. 
Then, combining \eqref{eq:br_policy} and  \eqref{eq:mixture},  the  fictitious play algorithm in \cite{heinrich2015fictitious}   computes a sequence of policies $\{ \pi_t \}_{t\geq 1}$.  
In particular, 
in the $t$-th iteration,  any agent $i$ first computes the $\epsilon_{t+1}$-best response policy $\tilde \pi_{t+1} ^i \in  \tt{Br}_{\epsilon_{t+1}}(\pi^{-i}_t)$ and then  constructs $\pi^i_{t+1}$ as  the mixture policy of $\pi^i_{t}$ and $\tilde \pi_{t+1}$ with weights $1- \alpha_{t+1}$ and $\alpha_{t+1}$, respectively.  Here,  both $\epsilon_t$  and $\alpha_t $ are taken to converge to zero as $t$ goes to infinity, and we further have $\sum_{t\geq 1} \alpha_t = \infty$. 
 We note, however, that although such a method provably converges to a Nash equilibrium of a zero-sum game via self-play, it suffers from the curse of dimensionality due to the need to iterate all states of the game in each iteration. For computational efficiency, \cite{heinrich2015fictitious} has also proposed a data-drive fictitious self-play framework where the best-response is computed via  fitted Q-iteration \cite{ernst2005tree, munos2007performance}  for the single-agent RL problem,  with the policy mixture being  learned through   supervised learning. This framework was later  adopted by  \cite{heinrich2014self, heinrich2016deep, kawamura2017neural,  zhang2019monte} to incorporate other single RL methods such as deep Q-network \cite{mnih2015human} and Monte-Carlo tree search  \cite{coulom2006efficient, kocsis2006bandit, browne2012survey}. Moreover, in a more recent work, \cite{perolat2018actor} has proposed a smooth fictitious  play algorithm \cite{fudenberg1995consistency} for zero-sum multi-stage games with simultaneous moves (a special case of zero-sum stochastic games). 
 Their algorithm combines the 
 actor-critic  framework \cite{konda2000actor}  with fictitious self-play, and infers the  opponent's policy implicitly via policy evaluation. Specifically, when the two players 
 adopt a joint policy $\pi = (\pi^1, \pi^2)$,  
 from the perspective of player $1$,  it infers $\pi^2$ implicitly by estimating $  \overbar Q _{\pi^1, \pi^2} $ via temporal-difference learning  \cite{sutton1987temporal}, where $ \overbar Q_{\pi^1, \pi^2} \colon \cS \times \cA^1 \rightarrow \RR$ is defined as 
 \$
 \overbar Q_{\pi^1, \pi^2} (s, a^1 ) = \EE\bigg[\sum_{t\geq 0}\gamma^t R^1(s_t,a_t,s_{t+1})\bigggiven s_0 = s, a_0^1 = a^1, a_0^2 \sim \pi^2 (\cdot \given  s) , a_t\sim \pi(\cdot\given s_t), \forall t\geq 1  \bigg] ,
 \$
 which is 
  the action-value function of player $1$ marginalized by $\pi^2$. Besides, the best response policy is obtained by taking the soft-greedy policy with respect to $\overbar Q_{\pi^1, \pi^2} $, i.e.,
  \#\label{eq:soft_greedy}
  \pi^1 (a^1  \given s)  \leftarrow \frac{\exp[  \eta ^{-1} \cdot \overbar Q_{\pi^1, \pi^2} (s,a^1 ) ] }{\sum_{a^1 \in \cA^1} \exp[  \eta^{-1} \cdot \overbar Q_{\pi^1, \pi^2} (s,a^1 ) ]  } ,
  \#
  where $\eta>0$ is the  smoothing parameter.  Finally, the algorithm is obtained by performing both policy evaluation and policy update in  \eqref{eq:soft_greedy} simultaneously using two-timescale updates \cite{borkar2008stochastic,kushner2003stochastic}, which ensure that the   policy updates, when using self-play, can be characterized by an ordinary differential equation whose asymptotically stable solution is 
  a smooth Nash equilibrium of the game.
 

 
 Another  family of popular  policy-based methods is based on the idea of \textit{counterfactural regret minimization} (CFR),   first proposed in \cite{zinkevich2008regret}, which has been a breakthrough in the effort to solve large-scale extensive-form games. 
Moreover, from a theoretical perspective, compared with fictitious play algorithms whose convergence is analyzed asymptotically via stochastic approximation, explicit regret bounds can be established using tools from online learning \cite{cesa2006prediction}, which yield rates of convergence to the  Nash equilibrium. 
  Specifically, 
  when $N$ agents play the extensive-form game for  $T$ rounds with  $\{ \pi_t \colon 1\leq t \leq T\}$, the \textit{regret} of player $i$   is defined~as 
  \#\label{eq:regret}
  \texttt{Reg}_T^i = \max_{\pi^i } \sum_{t = 1}^T \bigl [ R^i (\pi^i, \pi^{-i} _t  ) - R^i(\pi_t^i, \pi_t^{-i} ) \bigr ] ,
  \#
  where the maximum is taken over all possible policies of player $i$.
  In the following, before we define the notion of counterfactual regret, we first introduce a few notations. 
 Recall that we had defined the reach probability $\eta_{\pi} (h)$ in \eqref{eq:reach_prod}, which can be factorized into the product of each agent's contribution. That is, for each $i \in \cU \cup \{c\}$, we can group the probability   terms involving $\pi^{i}$ into $ \eta_{\pi}^i (h) $ and write 
 $
\eta_{\pi}(h) = \prod_{i \in \cN \cup \{c\} }  \eta_{\pi}^i (h) =  \eta_{\pi}^i (h) \cdot  \eta_{\pi}^{-i} (h).
$
Moreover, for any two histories $h , h' \in \cH$ satisfying $ h\sqsubseteq h'$, 
we define the \textit{conditional reach probability} 
$\eta_{\pi} ( h' \given h) $ as $ \eta _{\pi}(h') /  \eta_{\pi}(h)$ and define $\eta_{\pi}^i( h' \given h) $ similarly.
For any $s \in \cS^i $ and any $a \in \cA(s)$, we define $\cZ (s,a) = \{(h,z) \in \cH \times \cZ\given h \in s , ha \sqsubseteq z \}$, which contains all possible pairs of  history in   information state $s$ and    terminal history  after taking action $a$ at $s$.
Then, the \textit{counterfactual value function} is defined as 
\#
Q^i_{\texttt{CF}} (\pi, s, a)  & = \sum_{(h, z) \in \cZ (s,a) } \eta_{\pi}^{-i} (h) \cdot \eta_{\pi} ( z \given ha)  \cdot R ^{i} (z), \label{eq:cf_value_fun1}
 \#
which is the expected utility obtained by agent $i$ given that it has played to reached state~$s$.  
We also define  
$
 V^i_{\texttt{CF}} (\pi, s )  =  \sum_{a \in \cA(s) } Q^i_{\texttt{CF}} (\pi, s, a) \cdot \pi^i (a \given s)$.
  Then the difference between $Q^i_{\texttt{CF}} (\pi, s, a)  $ and $V^i_{\texttt{CF}} (\pi, s ) $ can be viewed as the value of action $a $ at information state $s \in \cS^i$, and  \textit{counterfactual regret} of agent $i$ at state $s$   is defined as 
\#\label{eq:immidiate_regret}
\texttt{Reg} _{T}^i (s ) = \max_{a \in \cA(s) } \sum_{i =1}^T \bigl [Q^i_{\texttt{CF}} (\pi_t, s, a)  - V^i_{\texttt{CF}} (\pi_t, s ) \bigr ], \qquad \forall s \in \cS^i. 
\#
Moreover, as shown in Theorem 3 of  \cite{zinkevich2008regret}, counterfactual regrets defined in \eqref{eq:immidiate_regret} provide an upper bound for the total regret in \eqref{eq:regret}: 
\#\label{eq:regret_upper}
\texttt{Reg}_T^i  \leq \sum_{s \in \cS^i } \texttt{Reg} _{T} ^{i , + }(s ) ,
\# 
where we let $x ^{+ }$ denote $\max \{ x, 0 \}$ for any $x \in \RR$. 
This bound lays the foundation of  \textit{counterfactual regret minimization} algorithms. Specifically, to minimize the total regret in \eqref{eq:regret}, it suffices to minimize the counterfactual regret for each information state, which can be obtained by any online learning algorithm, such as  EXP3   \cite{auer2002nonstochastic}, Hedge \cite{vovk1990aggregating, littlestone1994weighted,
 freund1999adaptive}, and regret matching \cite{hart2000simple}. %Exponential-weight algorithm for Exploration and Exploitation
All these methods ensure that the counterfactual regret is $   \cO(\sqrt{T})$ for all $s \in \cS^i$, which leads to an  $\cO( \sqrt{T})$ upper bound of the total regret. Thus,  applying CFR-type methods with self-play to a zero-sum two-play extensive-form game, the average policy is   an $\cO(\sqrt{1/T})$-approximate Nash equilibrium after $T$ steps.
In particular, 
 the vanilla  CFR algorithm updates the policies via regret matching, which yields that   
$  
 \texttt{Reg} _{T}^i (s )  \leq R_{\max}^i \cdot \sqrt{ A^i   \cdot T} 
$
for all $s\in \cS^i$, where we have introduced   
$$R_{\max}^i = \max_{z \in \cZ} R^i(z) - \min_{z \in \cZ} R^i(z), \qquad A_i = \max_{h \colon \tau(h) = i} | \cA(h ) |. $$
Thus, by  \eqref{eq:regret_upper},  the total regret of agent $i$ is bounded by $R_{\max}^i \cdot | \cS^i | \cdot  \sqrt{ A^i   \cdot T}$.  

One drawback of vanilla CFR is that the entire game tree needs to be traversed in each iteration, which can be computationally prohibitive. A number of  CFR variants have been proposed since   
the  pioneering work \cite{zinkevich2008regret} for  improving   computational efficiency. For example,  \cite{lanctot2009monte, burch2012efficient,gibson2012generalized, johanson2012efficient, lisy2015online,schmid2019variance} combine CFR with Monte-Carlo sampling;   \cite{waugh2015solving,morrill2016using, brown2019deep} propose to   estimate the counterfactual value functions via regression;  \cite{brown2015regret,brown2017dynamic, brown2017reduced}   improve the computational efficiency by  pruning  suboptimal paths in the game tree;    \cite{tammelin2014solving,tammelin2015solving, burch2019revisiting}  analyze the performance of a modification named  $\text{CFR}^{+}$, and \cite{zhou2018lazy} proposes lazy updates  with a near-optimal regret upper bound.   
 
 Furthermore, it has been  shown recently in \cite{srinivasan2018actor} that CFR is closely related to policy gradient methods. To see this, for any joint policy $\pi$ and any $i \in \cN$,
 we define the action-value function of agent $i$, denoted by $ Q_{\pi}^i $,  as 
 \#\label{eq:Q_efg}
 Q_{\pi}^i (s, a) =  \frac{1} {\eta_{\pi} (s) }  \cdot  \sum_{(h, z) \in \cZ (s,a) } \eta_{\pi}  (h)   \cdot \eta_{\pi} ( z \given ha)  \cdot R ^{i} (z),  \qquad \forall s \in \cS^i, \forall a \in \cA(s).
 \#
 That is,  $Q_{\pi}^i (s, a)$ 
 is the expected utility of agent~$i$ when the agents follow policy $\pi$ and agent $i$ takes action $a$ at information state $s \in \cS^i$, conditioning on $s $ being  reached. It has been  shown   in \cite{srinivasan2018actor} that the $Q^i_{\texttt{CF}}$  in 
 \eqref{eq:cf_value_fun1} is connected with $Q_{\pi}^i $ in \eqref{eq:Q_efg} via 
  $Q^i_{\texttt{CF}} (\pi, s, a)  =Q_{\pi}^i (s,a) \cdot  [ \sum_{h \in s } \eta^{-i} _{\pi} (h) ]   $. 
Moreover, in the tabular setting where we regard the joint policy $\pi$ as a table $\{ \pi^i (a\given s) \colon s \in \cS^i, a \in \cA(s) \}$, for any $s \in \cS^i$ and any $a \in \cA(s)$,  the policy gradient of $R^i (\pi) $ can be written~as 
  \$
\frac{\partial R^i (\pi) } { \partial _{\pi^i(a\given s) } }  =  \eta_{\pi}(s) \cdot  Q_{\pi}^i (s, a) = \eta_{\pi}^i (s) \cdot Q^i_{\texttt{CF}} (\pi, s, a) , \qquad \forall s \in \cS^i, \forall a \in \cA(s). 
  \$
  As a result, the advantage actor-critic  (A2C) algorithm \cite{konda2000actor}   is equivalent to a particular  CFR  algorithm, where the policy update rule is specified by the \textit{generalized infinitesimal gradient ascent} algorithm \cite{zinkevich2003online}.
 Thus,  \cite{srinivasan2018actor} proves that the regret of the  tabular A2C algorithm    is bounded by 
$
  | \cS^i | \cdot [  1 + A^i \cdot (R^i_{\max} )^2 ] \cdot  \sqrt{T}.     $    
  Following this work,  \cite{omidshafiei2019neural} shows that A2C where the policy is tabular and  is parametrized by a softmax function  is equivalent to CFR that uses  Hedge to update the policy. 
  Moreover,  \cite{lockhart2019computing} proposes a policy optimization method known as \textit{exploitability descent}, where the policy is updated using actor-critic, assuming the opponent plays the best counter-policy. This method is equivalent to the CFR-BR algorithm \cite{johanson2012finding} with Hedge.  
Thus,   \cite{srinivasan2018actor,omidshafiei2019neural, lockhart2019computing} show that actor-critic and policy gradient  methods for MARL can be formulated as  CFR methods and thus convergence to a Nash equilibrium of a zero-sum extensive-form game is guaranteed.
         
%  {\color{red} delete this part?
%  Neural replicator dynamics v.s. policy gradient, and connection to 
%evolutionary game theory. Other evolutionary  game theory perspective on MARL \cite{hofbauer2003evolutionary,bloembergen2015evolutionary}, and very recent empirical works that may be used for motivation \cite{khadka2018evolution,pmlr-v97-khadka19a,khadka2019evolutionary} }
%%\item Deep learning methods 
  
  
  In addition, besides   fictitious play and CFR methods introduced above,  multiple policy optimization methods have been  proposed for special classes of two-player zero-sum stochastic games  or extensive form games. For example, Monte-Carlo tree search methods
   have been  proposed  for perfect-information extensive games with simultaneous moves. 
It has been shown in \cite{schaeffer2009comparing} that  the MCTS  methods with UCB-type action selection rules, introduced in \S\ref{sec:value_comp}, fail to converge to a Nash equilibrium in simultaneous-move games, as UCB does not take into consideration the possibly adversarial moves of the opponent. 
To remedy this issue, 
   \cite{lanctot2013monte,lisy2013convergence, tak2014monte, kovavrik2018analysis} 
 have proposed to adopt stochastic policies and using  Hannan consistent methods such as EXP3   \cite{auer2002nonstochastic} and regret matching \cite{hart2000simple} to update the policies. With self-play, \cite{lisy2013convergence} shows that the average policy obtained by MCTS with any $\epsilon$-Hannan consistent policy update method  converges to an $\cO(D^2 \cdot \epsilon)$-Nash equilibrium, where $D$ is the maximal depth.  
 
 
 Finally, there are surging interests  in investigating   
 policy gradient-based methods in \emph{continuous} games, i.e., the games with continuous state-action spaces. With policy parameterization, finding the NE of  zero-sum Markov games becomes   a nonconvex-nonconcave saddle-point  problem in general \cite{mazumdar2018convergence,mazumdar2019finding,zhang2019policyb,bu2019global}. This hardness is inherent, even in the simplest linear quadratic setting with linear function approximation \cite{zhang2019policyb,bu2019global}. As a consequence,     most of the convergence results are \emph{local}  \cite{mescheder2017numerics,mazumdar2018convergence,adolphs2018local,daskalakis2018limit,mertikopoulos2019optimistic,fiez2019convergence,mazumdar2019finding,jin2019minmax}, in the sense that they address  the convergence behavior around  local NE points. Still, it has been shown that the vanilla gradient-descent-ascent (GDA) update, which is equivalent to the policy gradient update in MARL, fails to converge to local NEs, for either the non-convergent behaviors such as limit cycling  \cite{mescheder2017numerics,daskalakis2018limit,balduzzi2018mechanics,mertikopoulos2019optimistic}, or the existence of non-Nash stable limit points for the GDA dynamics \cite{adolphs2018local,mazumdar2019finding}. Consensus optimization \cite{mescheder2017numerics},   symplectic gradient adjustment \cite{balduzzi2018mechanics}, and extragradient method   \cite{mertikopoulos2019optimistic} have been advocated to mitigate the oscillatory behaviors around the equilibria; while \cite{adolphs2018local,mazumdar2019finding} exploit the curvature information so that all the stable limit points of the proposed updates are local NEs. 
Going beyond Nash equilibria, \cite{jin2019minmax,fiez2019convergence} consider gradient-based learning for \emph{Stackelberg equilibria}, which correspond to only the one-sided equilibrium solution in zero-sum games, i.e., either minimax or maximin, as the order of which player acts first is vital in nonconvex-nonconcave problems.  
\cite{jin2019minmax} introduces the  concept of \emph{local minimax} point as the solution, and shows that GDA converges to  local minimax points under mild conditions.  
\cite{fiez2019convergence}  proposes a two-timescale algorithm where the follower uses a gradient-play update rule, instead of an exact best response strategy, which has been shown to converge to the Stackelberg equilibria. 
 Under a stronger assumption of \emph{gradient dominance}, \cite{sanjabi2018solving,nouiehed2019solving} have shown that nested gradient descent methods converge to the stationary points of the outer-loop, i.e., minimax, problem at a sublinear rate. 
 
% has also be investigated for \emph{continuous} games, i.e., 
 
 
%when the state and action spaces are both continuous, 

 
%{\color{red} In addition, for two-player Markov games with linear-quadratic structures / continuous games (mention what it means by ``continuous'')


%\cite{mazumdar2018convergence} investigates the convergence of  gradient-based learning dynamics in continuous games, showing that vanilla policy gradient may converge to some non-Nash policies, even for zero-sum games.  



%  Policy gradient for zero-sum continuous games/saddle-point problems  \cite{mazumdar2019finding,jin2019minmax,zhang2019policyb}
 
 
 

We note that these convergence results have been developed for \emph{general}  continuous  games with \emph{agnostic} cost/reward functions, meaning that the functions may have various forms, so long as they are \emph{differentiable}, sometimes even \emph{(Lipschitz) smooth}, w.r.t.  each agent's policy parameter. For MARL, this is equivalent to requiring   differentiability/smoothness of the long-term \emph{return}, which relies on the properties of the  game, as well as of  the policy parameterization. Such an assumption is generally very restrictive. For example,   the Lipschitz smoothness assumption fails to hold globally for LQ games \cite{zhang2019policyb,mazumdar2019policy,bu2019global}, a special type of MGs.  Fortunately, thanks to the special structure of the LQ setting, \cite{zhang2019policyb} has proposed several   projected nested policy gradient methods that are guaranteed to have \emph{global} convergence to the NE, with convergence rates established.  This appears to be the first-of-its-kind result in MARL.  The results have then been improved by the techniques in a subsequent work of the authors  \cite{zhang2019policy}, which can  remove the projection step in the updates, for a more general class of such games. 
Very recently, \cite{bu2019global} also improves the results in  \cite{zhang2019policyb} independently, with different techniques.  

% }
   
        
  

\subsection{Mixed Setting}\label{subsec:mixed_setting}

In stark contrast with the fully  collaborative  and fully competitive settings, the mixed setting is notoriously challenging and thus  rather less well understood. 
Even in the simplest case of a  two-player general sum normal-form game, finding a  Nash  equilibrium is PPAD-complete \cite{chen2009settling}.   
Moreover, \cite{zinkevich2006cyclic} has proved  that value-iteration methods fail to find stationary Nash or correlated equilibria for general-sum Markov games.  Recently,   it is shown that vanilla  policy-gradient   methods avoid a non-negligible subset of   Nash equilibria in 
  general-sum continuous games \cite{mazumdar2018convergence},  including the  LQ general-sum games \cite{mazumdar2019policy}.   
Thus,  additional structures on either the games or the algorithms   
need to be exploited, to ascertain provably convergent MARL in the mixed setting. 


\vspace{6pt}
\noindent{\bf Value-Based Methods}
\vspace{3pt} 
 
Under relatively stringent assumptions, 
several value-based methods that extend 
Q-learning \cite{watkins1992q} 
 to the mixed setting are guaranteed to find an equilibrium.  
In particular, \cite{hu2003nash}  has proposed the Nash-Q learning  algorithm for general-sum Markov games, where  one maintains $N$  action-value functions $ Q_{\cN} = (Q^1, \ldots, Q^N) \colon    \cS \times \cA \rightarrow \RR^N  $  for all $N$ agents,  which are  updated 
 using sample-based estimator of a Bellman operator. Specifically,  letting $ R_{\cN} = ( R^1, \ldots, R^N)$ denote the reward functions of the agents,  Nash-Q uses the following Bellman operator: 
 \# \label{eq:nash_bellman}
   ( \cT^* Q_{\cN} )(s,a )  = \EE _{  s' \sim \cP(\cdot \given s, a )} \big  \{ R_{\cN}  (s, a , s' ) + \gamma \cdot \texttt{Nash}\bigl[  Q_{\cN} (s', \cdot ) \bigr ] \bigr \} , \quad \forall (s , a) \in \cS\times \cA,
   \#
where   $\texttt{Nash}[  Q_{\cN} (s', \cdot )   ]$  is the objective value of the Nash equilibrium of the stage game with rewards $\{ Q_{\cN} (s', a ) \}_{a \in \cA}   $. For zero-sum games, we have $Q^1 = -Q^2$ and thus the Bellman operator defined in \eqref{eq:nash_bellman} is equivalent to the one in  \eqref{eq:minimax_Q} used by minimax-Q learning \cite{littman1994markov}.  Moreover, \cite{hu2003nash} establishes  convergence to Nash equilibrium under the    restrictive assumption that  $\texttt{Nash}[  Q_{\cN} (s', \cdot )   ]$ in each iteration of the algorithm has unique Nash equilibrium.  
In addition, \cite{littman2001friend} has proposed the 
   Friend-or-Foe Q-learning   algorithm where each agent views the  other agent as either a ``friend'' or a ``foe''. In this case, $\texttt{Nash}[  Q_{\cN} (s', \cdot )   ]$ can be efficiently computed via linear programming.  This algorithm can be viewed as a  generalization of minimax-Q learning, and Nash equilibrium convergence is guaranteed for two-player zero-sum games and coordination games with a unique equilibrium. Furthermore,  
 \cite{greenwald2003correlated} has proposed  correlated Q-learning, which replaces $\texttt{Nash}[  Q_{\cN} (s', \cdot )   ]$ in \eqref{eq:nash_bellman} by computing a correlated equilibrium \cite{aumann1974subjectivity}, a more general equilibrium concept than Nash equilibrium. In a recent work, \cite{pmlr-v54-perolat17a}  has proposed a batch RL method to find an approximate  Nash equilibrium  via    Bellman residue minimization \cite{maillard2010finite}. They have proved that the global minimizer of the empirical Bellman residue is an approximate  Nash equilibrium, followed by the error propagation analysis for the algorithm.  
Also in the batch RL regime, \cite{zhang2018finite} has considered a simplified mixed setting for decentralized MARL: two teams of cooperative networked agents compete in a zero-sum Markov game. 
A decentralized  variant of FQI, where the agents within one team cooperate to solve \eqref{equ:FQI_coop} while the two teams essentially solve \eqref{eq:FQI}, is proposed. 
Finite-sample error bounds have  then been established for the proposed algorithm. 

%   \textcolor{red}{
%   More general function classes can be incorporated into minimax-Q learning 
%  under the framework of batch reinforcement learning \cite{lange2012batch}, where we iteratively update the action-value function by fitting   least-squares regression using  the target values  as the responses. 
%  Specifically, let $\cF$ be the function class of interest and let $  \{ (s_i, a_i^1, a_i^2, r_i, s_i' ) \}_{i\in[n]}$ be the dataset. For any $t \geq 0$,   let $Q_{t}$ be the current iterate in the $t$-th iteration. We define $y_i = r_i + \gamma \cdot \texttt{Value} [ Q_{t} (s_i', \cdot , \cdot )]$ for all $i \in [n]$. Then we update $ Q_{t}$   by solving a least-squares regression problem in $\cF$, that is, 
%  \#\label{eq:FQI}
%  Q_{t+1} \leftarrow \argmin_{ f\in \cF} \frac{1}{2n} \sum_{i=1}^n \bigl [ y_i - f(s_i, a_i^1, a_i^2 ) \bigr ]^2.
%  \#
%  This method generalizes the fitted Q-iteration (FQI) algorithm for MDPs to zero-sum games \cite{ernst2005tree, munos2007performance}. Moreover, when $\cF$ is the family of neural networks, the update in \eqref{eq:FQI} is the counterpart of the neural-FQI algorithm, which motivates the celebrated deep Q-network algorithm \cite{mnih2015human}.
%The  sample complexity of batch reinforcement learning methods for zero-sum Markov games can be obtained following the \textit{error propagation} framework proposed for MDPs 
%\cite{farahmand2010error, munos2007performance, antos2008learning,
%farahmand2016regularized}. Specifically,  \cite{yang2019theoretical} establish the sample complexity of the neural-FQI in \eqref{eq:FQI}, \cite{zhang2018finite}
%generalizes FQI to zero-sum Markov games between two teams over communication networks, \cite{pmlr-v37-perolat15,perolat2016use,perolat2016softened} study the sample complexity of approximate   policy iteration  methods \cite{scherrer2012approximate} for zero-sum games.}
%\issue{Need to echo back to \S\ref{subsubsec:networked_paradigm} and be simplified. Kaiqing's TO DO.}


To address  the scalability issue, independent learning is preferred, which, however, fails to converge in general  \cite{tan1993multi}.      \cite{arslan2017decentralized} has proposed  \emph{decentralized Q-learning}, a two
timescale modification of Q-learning, that is guaranteed to converge to the equilibrium  for \emph{weakly acyclic Markov games} almost surely. Each agent therein only  observes local action and reward, and neither observes nor  keeps track of others' actions.  
All agents are instructed to use the same  stationary \emph{baseline policy} for many consecutive stages, named \emph{exploration phase}. At the end of the \emph{exploration phase}, all agents are \emph{synchronized} to update their baseline policies, which makes the environment stationary for long enough, and enables the convergence of Q-learning based methods.  
Note that these algorithms can also be applied  to the cooperative setting, as these games include Markov teams as a special case. 

 

 
 \vspace{6pt}
\noindent{\bf Policy-Based Methods}
\vspace{3pt} 
  
For continuous games, 
due to the general negative  results  therein,  \cite{mazumdar2018convergence} introduces  a new class of games, \emph{Morse-Smale games}, for which the gradient dynamics correspond to gradient-like flows. Then, definitive statements on  almost sure convergence of PG methods  to either limit cycles, Nash equilibria, or non-Nash fixed points can be made, using tools from  dynamical systems theory.  
  Moreover,  \cite{balduzzi2018mechanics,letcher2019differentiable} have studied   the second-order structure of game dynamics, by decomposing  it into two components.The first one, named symmetric component,  relates to potential games, which yields gradient descent on some implicit function; the second one, named antisymmetric component, relates  to  \emph{Hamiltonian games} that follows some conservation law,  motivated by classical mechanical systems analysis. The fact that gradient descent converges to the Nash equilibrium of both types of games motivates the development of  the Symplectic Gradient Adjustment  algorithm that finds  \emph{stable fixed points} of the game, which constitute all local Nash equilibria for zero-sum games, and only 
  a subset of local NE for general-sum games. 
\cite{chasnov2019convergence} provides finite-time local convergence guarantees to  a neighborhood of a \emph{stable} local Nash equilibrium of continuous games, in both deterministic setting, with exact PG, and stochastic setting, with unbiased PG estimates. Additionally, \cite{chasnov2019convergence} has also explored the effects of \emph{non-uniform} learning rates on the learning dynamics and convergence rates. 
\cite{fiez2019convergence} has also considered \emph{general-sum} Stackelberg games, and shown that the same  two-timescale algorithm update as in the zero-sum case now converges almost surely to  the stable attractors only. It has also established  finite-time performance  for local convergence to a neighborhood of a stable Stackelberg equilibrium.  
In complete analogy to the zero-sum class, these convergence results for continuous games do not apply to MARL in Markov games directly, as they are  built upon the differentiability/smoothness of the long-term return, which may not hold for  general MGs, for example, LQ games \cite{mazumdar2019policy}.  Moreover, most of these convergence results are local instead of global. 
 
 
Other than continuous games,  
the policy-based methods summarized  in \S\ref{sec:policy_comp} can also be applied to the mixed setting via self-play. 
The validity of such an approach is based a fundamental connection between game theory and online learning -- If the external regret of each agent is no more than $\epsilon$, then their  average policies constitute an $\epsilon$-approximate
 \textit{coarse correlated equilibrium}  \cite{hart2000simple,hart2001general, hart2003uncoupled} of the general-sum normal-form games. Thus, although in general we are unable to find a Nash equilibrium, policy optimization with self-play   guarantees to  find a coarse correlated equilibrium in these normal-form games. 


\vspace{6pt}
\noindent{\bf Mean-Field Regime}
\vspace{3pt}

The scalability issue in the non-cooperative setting can also be alleviated in the mean-field  regime, as the cooperative setting discussed in \S\ref{subsubsec:homo_agents}. 
For general-sum games, 
\cite{pmlr-v80-yang18d} has proposed a modification of the Nash-Q learning algorithm where the actions of other agents are approximated by their empirical average. 
That is, the action value function of each agent $i$ is parametrized by $Q^i (s, a^i, \mu_{a^{-i}} )$, where $\mu_{a^{-i}}$ is the empirical distribution of  $\{a_j \colon j \neq i \}$.
 Asymptotic convergence of this mean-field  Nash-Q learning algorithm has also been established.
 
 
Besides, most mean-field RL algorithms are focused on addressing the   \textit{mean-field game} model.  
%MFG is 
%first introduced independently by  \cite{huang2003individual,huang2006large,lasry2007mean,bensoussan2013mean} for modeling games where a large number of indistinguishable agents interact with each other and with the stochastic environment, with each agent having infinitesimal effect on  the system.
In mean-field games, each agent $i$ has a local state $s^i \in \cS $ and a local action $a^i \in \cA$, 
and  the interaction among other agents is  captured  by an aggregated effect $\mu$, also known as the \emph{mean-field term}, which is a functional of the empirical distribution of the local states and actions of the agents.  
Specifically, at the $t$-th time step, when  agent $i$ takes action $a_t^i$ at state $s_t^i$ and the mean-field term is $\mu_t$, it receives an immediate reward $R(s_t^i, a_t^i, \mu_t)$ and its local state evolves into $  s_{t+1}^i \sim \cP(\cdot \given s_t^i, a_t^i, \mu_t) \in \Delta(\cS)$. Thus, from the perspective of agent $i$,  instead of participating in a multi-agent game, it is faced with a time-varying MDP parameterized by the sequence of mean-field terms $\{ \mu_t \}_{t\geq 0}$, which in turn is determined by the states and actions of all agents. 
The solution concept in MFGs is the   \emph{mean-field equilibrium}, which is a sequence of \emph{pairs} of policy and mean-field terms
$\{\pi^*_t, \mu^*_t \}_{t\geq 0}$ that satisfy the following two conditions: (1) $\pi^* = \{\pi^*_{t} \}_{t\geq 0}$ is the optimal policy for the time-varying MDP specified by $\mu^* = \{\mu^*_t \}_{t\geq 0}$, and (2)  $\mu^*$ is generated when each agent follows policy $\pi^*$.  The existence of the mean-field equilibrium for discrete-time MFGs has been  studied in  
\cite{saldi2018markov, saldi2019approximate,saldi2019discrete,saldi2018discrete}
and their constructive  proofs exhibit that the mean-field equilibrium can be obtained via a fixed-point iteration. Specifically, one can construct a sequence of policies and mean-field terms $\{\pi^{(i)} \}_{i\geq 0}$ and $\{ \mu^{(i)}\} _{i\geq 0}$ such that $\{ \pi^{(i)}\}_{t\geq 0}$ solves the time-varying MDP specified by $\mu^{(i)}
$, and $\mu^{(i+1)}$ is generated when all players adopt policy $\pi^{(i)}$. 
Following this agenda, various model-free RL methods are proposed for solving MFGs where $\{\pi^{(i)} \}_{i\geq 0}$ is approximately solved via single-agent RL such as Q-learning \cite{guo2019learning} and policy-based methods 
\cite{subramanian2019reinforcement, fu2019actor}, with  $\{ \mu^{(i)}\} _{i\geq 0}$ being  estimated via sampling.  
In addition,   \cite{hadikhanloo2019finite,elie2019approximate} recently propose   fictitious play updates for the mean-field state where we have $\mu^{(i+1)} = (1-\alpha^{(i)}) \cdot \mu^{(i)} + \alpha^{(i)}  \cdot \hat  \mu^{(i+1)} $, with $ \alpha^{(i)}$
 being the learning rate and $\hat  \mu^{(i+1)}$ being the mean-field term generated by policy $\pi^{(i)}$. Note that the aforementioned works focus on  the settings with either \emph{finite} horizon \cite{hadikhanloo2019finite,elie2019approximate} or \emph{stationary} mean-field equilibria  \cite{guo2019learning,subramanian2019reinforcement,fu2019actor}  only. Instead, recent works 
 \cite{anahtarci2019value,zaman2019approx}  consider possibly non-stationary mean-field equilibrium in infinite-horizon settings, and develop equilibrium computation algorithms that have laid  foundations for model-free RL algorithms. 
 
 




 



\iffalse
 
%\begin{comment}
\vspace{10pt}
\noindent{XXXXXXXXXXXXX}
\issue{@Zhuoran: Plz check out and clean the comments you left.}
XXXXXXXXXXXXXX

\noindent\issue{What I haven't coverred}

\begin{itemize}
	%\item Nash Q learning \cite{hu2003nash}, Friend-or-Foe Q-learning \cite{littman2001friend}, and correlated-Q learning \cite{greenwald2003correlated} that generalizes the two (convergence to correlated equilibrium).  
	%\item {Zinkevich's negative result \cite{zinkevich2006cyclic}: \emph{value-based} methods are not sufficient for general-sum Markov games.}
%\item Policy gradient for general-sum continuous games \cite{mazumdar2018convergence}
% \item  Batch ADP papers for general sum Markov games \cite{pmlr-v54-perolat17a} 
%\item ICML differentiable games paper (best-paper runner-up) for continuous games \cite{balduzzi2018mechanics}, and longer version \cite{letcher2019differentiable}.
%\item Neural replicator dynamics v.s. policy gradient, and connection to evolutionary game theory \cite{omidshafiei2019neural}. Other evolutionary  game theory perspective on MARL \cite{hofbauer2003evolutionary,bloembergen2015evolutionary}, and very recent empirical works that may be used for motivation \cite{khadka2018evolution,pmlr-v97-khadka19a,khadka2019evolutionary}
%\item Deep learning methods 
%\item Mean-field approximation: (i) mean-field action approximation   \cite{pmlr-v80-yang18d} (it indeed has theory); (ii) standard  MFG-based: Mahajan's policy gradient \cite{subramanian2019reinforcement}; Guo's Q-learning \cite{guo2019learning};  fictitious play for (continuous-time) potential MFG \cite{cardaliaguet2017learning}, and (discrete-time) but finite-time and finite-state MFG \cite{hadikhanloo2019finite}; approximate fictitious play for MFG \cite{elie2019approximate}; actor-critic for LQ-MFG [zhuoran's paper], Q-learning for LQ-MFG with non-stationary MF [Aneeq's paper]; and references therein.

 
\item RL in (LQ) games over graphs/graphical games: discrete-time \cite{abouheaf2014multi} and  continuous-time \cite{vamvoudakis2012multi}; off-policy \cite{li2017off}
\item A survey from control on MARL \cite{vamvoudakis2017game} (mostly continuous-time settings) 

\item Add some new papers from the CMU group on EFG. These methods use the sequence form. Also cite deepmind's new nature paper.

\end{itemize}

\noindent{XXXXXXXXXXXXXXXXXXXXXXXXXXXXXXXXXXXXXXXXXXXXXXXXXXXXXXXXXXXXXXXXX}
%\end{comment}

 \fi 



